\begin{abstract}
Existing financial multi-agent systems rely on verbal reinforcement for self-improvement, rewriting reflections into prompt prefixes without updating model parameters or long-term memory. This approach fails under regime shifts: live trading benchmarks report net losses for frontier models, and reinforcement-finetuned agents achieve higher directional accuracy yet lower cumulative returns. We propose \textsc{FinOPD}, a framework that enables policy-level self-evolution in financial multi-agent systems through three mechanisms: (1) on-policy self-distillation conditioned on hindsight risk-adjusted PnL as privileged information, which anchors learning to profitability rather than prediction accuracy; (2) Shapley-weighted distillation for multi-agent credit assignment; and (3) a non-parametric belief store that preserves winning decision patterns alongside parametric adaptation. On Dow-30 and CSI-300 benchmarks, \textsc{FinOPD} obtains state-of-the-art Sharpe ratio, maximum drawdown, and cross-regime robustness over twelve baselines, with sustained improvement across self-evolution iterations under strict post-cutoff evaluation.
\end{abstract}
% 中文翻译：现有金融多智能体系统仅依赖语言强化进行自我改进，将反思重写为 prompt 前缀而不更新模型参数或长期记忆。这种方法在市场 regime 转换时失效：实盘基准显示前沿模型出现净亏损，强化微调的 agent 方向准确率提高但累计收益下降。我们提出 FinOPD，通过三种机制实现金融多智能体系统的策略级自进化：(1) 以事后风险调整 PnL 为特权信息的在线策略自蒸馏，将学习锚定于盈利性而非预测准确率；(2) 基于 Shapley 值加权的蒸馏解决多智能体信用分配；(3) 非参数 belief 存储在参数适应之外保留获胜决策模式。在 Dow-30 和 CSI-300 基准上，FinOPD 在夏普比率、最大回撤和跨 regime 鲁棒性上超越十二个基线，在严格的 post-cutoff 评估下自进化迭代中持续改进。
